% introduction.tex -- Section 1: Introduction

\section{Introduction}
\label{sec:introduction}

Vision token pruning is one of the most reliable efficiency techniques in Vision Transformer (ViT) inference.
DynamicViT~\cite{rao2021dynamicvit} removes 66\% of patch tokens with less than 0.5\% accuracy loss on ImageNet classification.
Token Merging (ToMe)~\cite{bolya2023tome} merges 40\% of tokens while retaining over 99\% of classification accuracy.
EViT~\cite{liang2022evit} discards 50\% of inattentive tokens with negligible degradation.
These consistent results across architectures and benchmarks have established a strong prior: vision tokens are highly redundant, and aggressive reduction is safe.
A natural question follows.
Vision-Language-Action (VLA) models~\cite{brohan2023rt2,kim2024openvla,black2024pi0,shukor2025smolvla} use the same ViT encoders as their visual front-ends, producing hundreds of patch tokens that dominate the input sequence to the language model backbone.
For real-time robotic control, where inference latency directly constrains action frequency, reducing these vision tokens would yield immediate speedups in the most expensive component of the forward pass.
Does ViT token reduction transfer to VLA closed-loop control?

This question has practical urgency.
Current VLA models such as RT-2~\cite{brohan2023rt2} and $\pi_0$~\cite{black2024pi0} use large vision encoders that produce 256 or more tokens per image.
When multiple camera views are used, the vision token count grows proportionally.
In a typical deployment scenario, a VLA must produce actions at 5--10 Hz to maintain smooth control, and the LLM backbone's self-attention over vision tokens is the dominant latency bottleneck because attention cost scales quadratically with sequence length.
SmolVLA~\cite{shukor2025smolvla}, a compact 451.5M-parameter VLA, produces 256 vision tokens from two cameras (128 per view), and these tokens constitute the majority of the LLM backbone's input sequence.
Reducing the vision token count from 256 to 128 would cut attention FLOPs by approximately 4$\times$ for the vision portion of the sequence, a significant gain for real-time control.
Several recent works have begun exploring token efficiency for Vision-Language Models (VLMs) and VLAs, motivated by these potential speedups.
FastV~\cite{chen2024fastv} identifies and prunes inattentive visual tokens in VLM inference, achieving 2$\times$ speedup with minimal quality loss on VLM benchmarks.
VLA-Cache~\cite{xu2025vlacache} caches vision tokens across timesteps to reduce redundant computation in temporal sequences.
VLA-IAP~\cite{cheng2026vlaiap} proposes importance-aware pruning for VLA inference acceleration.
These proposals are motivated by the same intuition that has driven ViT pruning: if many vision tokens contribute little to the output, removing them should provide efficiency gains at minimal performance cost.
The implicit assumption is that the redundancy structure observed in ViT classification (where attention concentrates on a few discriminative patches) also holds in VLA models.

Despite this growing interest, no prior work has systematically evaluated inference-time vision token reduction in VLA closed-loop control.
Existing evaluations of VLM token reduction rely on offline metrics such as action loss, perplexity, or next-token prediction accuracy~\cite{chen2024fastv,cheng2026vlaiap}.
These proxies measure average-case behavior over individual timesteps but do not capture the compounding effects of action errors in closed-loop execution, where a single bad action can move the robot into an unrecoverable state.
The distinction between open-loop and closed-loop evaluation is well-established in robotics: policies that achieve low imitation loss can still fail catastrophically when errors compound across timesteps in the absence of corrective demonstrations.
For token reduction specifically, even a small perturbation to the action output at one timestep can push the robot off the demonstrated trajectory, and subsequent observations will differ from the training distribution, triggering further errors in a cascading failure.
The robotics community currently lacks empirical guidance on whether vision token reduction is a viable efficiency strategy for deployed VLA policies, or whether offline metrics provide misleading confidence about closed-loop safety.
This gap is consequential: if practitioners apply ViT pruning techniques to VLA deployment based on encouraging offline numbers alone, they risk deploying policies that fail silently in the real world.

We conduct the first systematic closed-loop evaluation of this question.
Using SmolVLA on the LIBERO benchmark~\cite{liu2023libero}, we test six pruning configurations spanning three axes: scoring function (a trained Attention-Conditioned Importance Scorer versus random selection), retention ratio ($k \in \{0.5, 0.7, 0.9\}$), and mechanism (zero-masking, true removal, and token merging).
The result is unambiguous.
Every single configuration collapses to 0\% Success Rate (SR) across 500 evaluated episodes (Wilson 95\% CI: [0.0\%, 0.76\%]), compared to a 44.8\% baseline SR achieved by the identical model checkpoint with pruning disabled.
There is no gradual degradation curve.
Even retaining 90\% of tokens, removing just 10\% of the vision input, produces complete task failure.
We term this the \emph{cliff edge}: a binary transition from functional to non-functional behavior with no intermediate regime.
The cliff edge is method-agnostic, appearing identically under learned scoring, random scoring, different retention ratios, and all three reduction mechanisms.
Cross-suite evaluation on LIBERO-Spatial (baseline 51.8\% SR) and LIBERO-Goal (baseline 63.0\% SR) confirms the same collapse.

We provide a mechanistic explanation for why VLA control exhibits this cliff edge while ViT classification degrades gracefully.
The explanation rests on three observations.
First, the LLM backbone in SmolVLA attends to vision tokens near-uniformly: the attention entropy ratio reaches 0.98 (6.86 out of 7.0 bits), meaning no subset of tokens is attended to preferentially.
This uniformity arises because SigLIP~\cite{zhai2023siglip} encodes non-overlapping spatial patches without a \texttt{[CLS]} aggregation token; each token carries unique, irreplaceable spatial information about the scene.
Second, removing any subset of tokens causes immediate catastrophic action disruption at the single-step level, with a relative action error of 111.5\% and gripper dimension error of 1.02 (complete open/close inversion).
The failure is not gradual error accumulation over time but instantaneous corruption of the action output.
Third, standard offline evaluation metrics (importance margins, averaged per-step action losses) provide misleadingly positive signals that fail to predict this closed-loop collapse.
The gap between offline proxies and closed-loop outcomes demonstrates that offline evaluation alone is insufficient for validating VLA compression strategies.
This mechanistic picture explains the fundamental difference between the two settings.
In ViT classification, the task requires mapping an image to a single discrete label, and the classifier can derive this label from a concentrated subset of discriminative patches; removing the remaining low-attention tokens preserves the decision boundary.
In VLA control, the policy must reconstruct the full spatial layout of the scene (object positions, gripper pose, spatial relationships between entities) to generate precise 7-DoF continuous actions, and the near-uniform attention distribution reflects this holistic spatial requirement.
The two settings have fundamentally different information requirements, and techniques designed for one do not transfer to the other.
We believe this negative result~\cite{karl2024negative,zeng2023transformers_tsf} carries important implications for the VLA efficiency community: future work on VLA compression should seek alternatives to naive token removal, such as training-time distillation into fewer tokens or architectural modifications that build in spatial redundancy.

Our contributions are:
\begin{enumerate}[leftmargin=*]
    \item \textbf{Cliff-edge phenomenon} (Section~\ref{sec:results}): We present the first systematic closed-loop evaluation of VLA vision token reduction, revealing that all tested configurations (6 conditions across 3 axes) collapse from 44.8\% SR to 0\% SR with no gradual degradation, a result that contradicts the transferability assumption from ViT pruning literature.
    \item \textbf{Mechanistic explanation} (Section~\ref{sec:analysis}): We identify the causal mechanism: near-uniform attention (entropy ratio 0.98) combined with instantaneous single-step action disruption (111.5\% relative error, gripper inversion) explains why VLA control exhibits a binary failure mode rather than graceful degradation.
    \item \textbf{Offline metric disconnect} (Section~\ref{sec:metric_disconnect}): We demonstrate that standard offline evaluation proxies (importance margins, per-step action losses) provide misleadingly optimistic estimates, establishing that closed-loop evaluation is essential for validating any VLA compression technique.
\end{enumerate}
